\chapter{El fenómeno: EEG y crisis epilépticas}
\label{cap:fenomeno}

Este capítulo describe qué se mide, por qué una crisis deja una huella en la señal y cómo se traduce
esa huella en un modelo matemático. El modelo de la ecuación~\eqref{eq:modelo_eeg} es la base de todo
lo que sigue: justifica el uso de la transformada de Fourier (Capítulo~\ref{cap:fourier}), los atributos de
potencia por banda, la calibración por paciente y el planteo de la detección como un test de hipótesis.

\section{Epilepsia y crisis}

Una \emph{crisis epiléptica} es la aparición transitoria de signos o síntomas debidos a una actividad
neuronal anormalmente excesiva o sincrónica en el cerebro. La \emph{epilepsia} es la enfermedad que
predispone a sufrir crisis en forma recurrente; en la práctica clínica se diagnostica, entre otros
criterios, tras dos crisis no provocadas separadas por más de 24 horas \parencite{fisher2014}.

Para el diagnóstico y el ajuste del tratamiento se registra el electroencefalograma (EEG) durante
horas o días. Las crisis ocupan una fracción muy chica de ese tiempo: en los seis pacientes de
CHB-MIT usados en este trabajo, 34 crisis suman el 0,37~\% de 222,8 horas de registro. Revisar todo
el registro a mano es lento, y esa es la motivación de la detección automática.

\section{Origen de la señal de EEG}

El EEG mide diferencias de potencial eléctrico generadas por la actividad del cerebro. La mayor
parte de la señal proviene de las corrientes postsinápticas de las neuronas piramidales de la corteza.
Estas neuronas están alineadas en forma perpendicular a la superficie cortical, de modo que cuando
miles de ellas reciben entradas sincronizadas sus corrientes se suman y el conjunto se comporta como
un \emph{dipolo} de corriente. Un potencial de acción aislado es demasiado breve y local para verse en
la superficie; lo que se registra es la suma de actividad lenta y sincronizada de millones de neuronas.

El potencial que genera ese dipolo atraviesa el líquido cefalorraquídeo, el cráneo y el cuero
cabelludo, que actúan como un conductor de volumen: atenúan la señal y la difunden espacialmente.
Por eso la amplitud típica del EEG de superficie es de decenas de microvoltios, mientras que el EEG
intracraneal (electrodos sobre la corteza o dentro de ella, como en los sets C, D y E de Bonn) alcanza
cientos.

La consecuencia importante para el modelado es que la amplitud registrada depende mucho de la
\emph{sincronía}: $M$ fuentes con fases aleatorias producen una amplitud que crece como $\sqrt{M}$,
mientras que $M$ fuentes en fase producen una amplitud que crece como $M$. Una crisis, que es por
definición actividad excesivamente sincrónica, se ve como un aumento grande de amplitud.

\section{Ritmos y bandas de frecuencia}

El EEG normal no es ruido blanco: tiene oscilaciones características (ritmos) que se describen por su
banda de frecuencia (Tabla~\ref{tab:bandas}).

\begin{table}[ht]
\centering
\caption{Bandas de frecuencia clásicas del EEG, con los límites usados en este trabajo.}
\label{tab:bandas}
\begin{tabular}{llp{8cm}}
\toprule
Banda & Rango [Hz] & Contexto típico \\
\midrule
delta $\delta$ & 0,5--4 & sueño profundo; actividad lenta patológica en vigilia \\
theta $\theta$ & 4--8 & somnolencia; frecuente en descargas ictales \\
alfa $\alpha$ & 8--13 & vigilia relajada con ojos cerrados, en zonas occipitales \\
beta $\beta$ & 13--30 & vigilia activa, concentración \\
gamma $\gamma$ & 30--40 & actividad rápida (acotada a 40 Hz por el filtro) \\
\bottomrule
\end{tabular}
\end{table}

Un ejemplo de los datos: en Bonn, el set B (voluntarios sanos, ojos cerrados) tiene el 47,6~\% de su
potencia en la banda alfa, contra el 24,1~\% del set A (los mismos voluntarios con los ojos abiertos).
Es el ritmo alfa occipital, que se atenúa al abrir los ojos (Capítulo~\ref{cap:fourier}).

\section{Qué cambia durante una crisis}

Durante una crisis aparece una \emph{descarga rítmica}: una oscilación de gran amplitud, casi
periódica, que puede empezar en una zona (crisis focal) y propagarse, o comprometer los dos hemisferios
desde el inicio (crisis generalizada). En los datos de este trabajo se observa que:

\begin{itemize}
  \item En Bonn, los segmentos ictales (set E) tienen un desvío estándar de 341~µV contra 48--90~µV
        de los demás sets, y entre 10 y 100 veces más potencia entre 3 y 25 Hz.
  \item En CHB-MIT la potencia media sube en \emph{todas} las bandas y en \emph{todos} los pacientes, entre 2 y
        50 veces; pero la banda que más sube depende del paciente: delta en chb01 y chb08, alfa y
        beta en chb05 y chb10 (Capítulo~\ref{cap:datos}).
  \item Después de la crisis queda un período de actividad lenta de gran amplitud (estado
        \emph{post-ictal}) que no se anota como crisis.
\end{itemize}

\section{Modelo matemático de la señal}
\label{sec:modelo}

\subsection{Montaje bipolar}

Cada electrodo $e$ registra un potencial $v_e(t)$ respecto de una referencia. En el montaje bipolar de
CHB-MIT, cada canal es la diferencia entre dos electrodos vecinos del sistema 10-20:
\begin{equation}
  x_{A\text{-}B}(t) = v_A(t) - v_B(t).
  \label{eq:bipolar}
\end{equation}
La resta cancela lo que es común a los dos electrodos (la referencia, buena parte de la interferencia
de red y de la actividad lejana) y resalta la actividad local entre ellos. Se usan los 18 canales de la
``doble banana''; en adelante $x_c(t)$, $c = 1,\dots,C$, denota el canal $c$.

\subsection{Modelo aditivo}

Se modela la señal registrada en el canal $c$ como la suma de cuatro componentes:
\begin{equation}
\boxed{\;
  x_c(t) \;=\; g_{p,c}\Big[\, \underbrace{b_c(t)}_{\text{fondo}}
  \;+\; \underbrace{\mathds{1}_{\mathcal{C}}(t)\, s_c(t)}_{\text{descarga ictal}} \,\Big]
  \;+\; \underbrace{a_c(t)}_{\text{artefactos}}
  \;+\; \underbrace{\eta_c(t)}_{\text{ruido}}
\;}
\label{eq:modelo_eeg}
\end{equation}
donde:
\begin{itemize}
  \item $b_c(t)$ es la \textbf{actividad de fondo}, modelada como una suma de osciladores en las bandas
        de la Tabla~\ref{tab:bandas} con amplitudes y fases que varían lentamente:
        \begin{equation}
          b_c(t) = \sum_{k} A_{c,k}(t) \cos\!\big(2\pi f_k t + \varphi_{c,k}(t)\big).
          \label{eq:fondo}
        \end{equation}
        En ventanas de pocos segundos se la trata como un proceso estacionario de media cero con densidad
        espectral $S_{b,c}(f)$.
  \item $\mathcal{C} = \bigcup_j [a_j, b_j]$ es el conjunto de intervalos con crisis, y
        $\mathds{1}_{\mathcal{C}}(t)$ vale 1 dentro de una crisis y 0 fuera.
  \item $s_c(t)$ es la \textbf{descarga ictal}, una oscilación casi periódica con frecuencia fundamental
        $f_0(t)$ (típicamente entre 3 y 8 Hz, que puede variar a lo largo de la crisis) y armónicos:
        \begin{equation}
          s_c(t) = \sum_{h=1}^{H} B_{c,h}(t) \cos\!\big(2\pi h\, f_0(t)\, t + \psi_{c,h}\big).
          \label{eq:descarga}
        \end{equation}
        Las amplitudes $B_{c,h}$ son grandes en los canales cercanos al foco y pueden ser casi nulas en los
        demás.
  \item $g_{p,c} > 0$ es una \textbf{ganancia} que depende del paciente $p$ y del canal: espesor del cráneo,
        impedancia de los electrodos, edad. Es la que hace que el ``nivel basal'' sea distinto entre
        pacientes.
  \item $a_c(t)$ son \textbf{artefactos} (movimiento, músculo, parpadeo, electrodo suelto) y $\eta_c(t)$ es
        ruido del equipo e interferencia de banda angosta (50 Hz en Bonn; 16 y 32 Hz en algunos registros
        de CHB-MIT).
\end{itemize}

\subsection{Consecuencia en el dominio de la frecuencia}

Si las componentes son procesos independientes y de media cero, sus densidades espectrales se
suman (el espectro de una suma de procesos no correlacionados es la suma de los espectros):
\begin{equation}
  S_{x_c}(f) = g_{p,c}^2 \Big[ S_{b,c}(f) + \mathds{1}_{\mathcal{C}}\, S_{s,c}(f) \Big] + S_{a,c}(f) + S_{\eta,c}(f).
  \label{eq:espectro_modelo}
\end{equation}
La potencia en una banda $\mathcal{B} = [f_1, f_2)$ es la integral de la densidad espectral,
$P_{c,\mathcal{B}} = \int_{\mathcal{B}} S_{x_c}(f)\,df$. Ignorando artefactos y ruido, el cociente entre la potencia
durante y fuera de la crisis es
\begin{equation}
  \frac{P^{\text{crisis}}_{c,\mathcal{B}}}{P^{\text{basal}}_{c,\mathcal{B}}}
  = 1 + \frac{\int_{\mathcal{B}} S_{s,c}(f)\,df}{\int_{\mathcal{B}} S_{b,c}(f)\,df} \;>\; 1,
\end{equation}
y en escala logarítmica
\begin{equation}
  \log_{10} P^{\text{crisis}}_{c,\mathcal{B}} - \log_{10} P^{\text{basal}}_{c,\mathcal{B}}
  = \log_{10}\!\left(1 + \frac{P_{s,c,\mathcal{B}}}{P_{b,c,\mathcal{B}}}\right) > 0.
  \label{eq:dif_log}
\end{equation}
Tres consecuencias directas del modelo:
\begin{enumerate}
  \item \textbf{La potencia por banda es un buen atributo.} La crisis suma potencia, sobre todo en las
        bandas que contienen $f_0$ y sus armónicos. Por eso la transformada de Fourier es la herramienta
        natural (Capítulo~\ref{cap:fourier}).
  \item \textbf{El logaritmo convierte la ganancia en una constante aditiva.} Como
        $\log P_{c,\mathcal{B}} = 2\log g_{p,c} + \log\!\big(P_{b,c,\mathcal{B}} + \mathds{1}_{\mathcal{C}}P_{s,c,\mathcal{B}}\big)$,
        el término $2\log g_{p,c}$ no depende del tiempo ni de la clase. Restar la mediana de cada
        paciente (que, al ser las crisis menos del 1~\% del tiempo, es esencialmente el nivel basal) lo
        elimina. Esta es la justificación de la \emph{calibración por paciente} que se usa en los
        Capítulos~\ref{cap:convexos} y~\ref{cap:cnn}.
  \item \textbf{La banda que más cambia depende de $f_0$ y de $S_{b,c}$}, que varían entre pacientes. Un
        detector que funcione para pacientes nuevos tiene que combinar varias bandas y canales.
\end{enumerate}

\section{La detección como test de hipótesis}
\label{sec:hipotesis}

Se divide la señal en ventanas de $L = 4$ s. Para cada ventana $w$, con sus datos $\mathbf{X}_w$ (los
18 canales), hay dos hipótesis:
\begin{equation}
  H_0:\ \mathbf{X}_w = \text{fondo} + \text{artefactos} + \text{ruido} \qquad\text{(sin crisis)},
  \qquad
  H_1:\ \mathbf{X}_w \text{ incluye además } s_c(t) \qquad\text{(crisis)}.
\end{equation}
Un detector calcula un puntaje $s(\mathbf{x})$ a partir de los atributos $\mathbf{x}$ de la ventana y
decide $H_1$ si $s(\mathbf{x}) > \tau$. El lema de Neyman-Pearson (Sección~\ref{sec:neyman_pearson})
dice que el mejor puntaje posible es la razón de verosimilitud $p(\mathbf{x}\mid H_1)/p(\mathbf{x}\mid H_0)$.
Como esas densidades son desconocidas, se aproximan con modelos entrenados con datos: regresión logística,
SVM y una red convolucional.

\paragraph{El problema del desbalance.} Con probabilidad a priori de crisis $\pi_1 \approx 0{,}0037$,
la regla de Bayes dice que la probabilidad de crisis dada una ventana es
\begin{equation}
  P(H_1 \mid \mathbf{x}) = \frac{\pi_1\, p(\mathbf{x}\mid H_1)}{\pi_1\, p(\mathbf{x}\mid H_1) + \pi_0\, p(\mathbf{x}\mid H_0)}.
\end{equation}
Para que esta probabilidad supere 1/2 hace falta que la razón de verosimilitud supere
$\pi_0/\pi_1 \approx 270$. Un clasificador que nunca detecta nada acierta el 99,6~\% de las ventanas; por
eso la exactitud (\emph{accuracy}) no sirve como métrica y en el Capítulo~\ref{cap:metricas} se usan
sensibilidad, especificidad, falsas alarmas por hora y curvas ROC.
